\section{Finite-set laws and terminal defects}

Define
\begin{equation}\label{eq:F-and-q}
  F(t)=2\Phi(\sqrt t/2)-1,
  \qquad
  q(t)=1-F(t)=2\Phi(-\sqrt t/2).
\end{equation}

\begin{corollary}[Exact finite-set stabilization]\label{cor:finite-set}
For a finite $A\subseteq I$ with $|A|=M$,
\begin{equation}\label{eq:max-cdf}
  \P(T_A\le t\mid X)=\P(T_A\le t)=F(t)^M,
  \qquad t\ge0.
\end{equation}
The one-coordinate density and Laplace transform are
\begin{equation}\label{eq:scalar-density-laplace}
  f_T(t)=\frac{e^{-t/8}}{\sqrt{8\pi t}},\quad t>0,
  \qquad
  \E e^{-sT}=(1+8s)^{-1/2},\quad s\ge0.
\end{equation}
Consequently,
\begin{equation}\label{eq:scalar-moments}
  \E T=4,
  \qquad
  \Var(T)=32.
\end{equation}
The maximum $T_A$ has density
\begin{equation}\label{eq:max-density}
  f_{T_A}(t)=M F(t)^{M-1}\frac{e^{-t/8}}{\sqrt{8\pi t}},
  \qquad t>0.
\end{equation}
\end{corollary}

\begin{proof}
Equation \eqref{eq:max-cdf} follows from Theorem~\ref{thm:main}.  The remaining
formulas follow by differentiating $F$ and using the standard gamma transform.
\end{proof}

At time $t$, call coordinate $i$ a \emph{remaining terminal defect} if it will
be decoded incorrectly at some time strictly greater than $t$.  Equivalently,
it is a defect when $T_i>t$.

\begin{corollary}[Exact defect count]\label{cor:defect-count}
For finite $A\subseteq I$, define
\begin{equation}\label{eq:defect-count}
  N_A(t)=\sum_{i\in A}\1\{T_i>t\}.
\end{equation}
Conditionally on $X$, and hence also unconditionally,
\begin{equation}\label{eq:binomial-defects}
  N_A(t)\sim\operatorname{Bin}(|A|,q(t)).
\end{equation}
In contrast, the number of coordinates decoded incorrectly at the single
snapshot time $t>0$ is binomial with parameter $q(t)/2$.
\end{corollary}

\begin{proof}
The first assertion follows from the independent marking theorem.  For a
snapshot error, either target bit requires a centered normal fluctuation of
magnitude at least $t/2$, so the conditional error probability is
$\Phi(-\sqrt t/2)=q(t)/2$.  Conditional independence follows from the
independence of the Brownian coordinates.
\end{proof}

Thus the factor of two in \eqref{eq:factor-two-intro} is exact: among the paths
that will err after time $t$, half are already in error at time $t$, while half
are currently correct and later return across the threshold.

\begin{corollary}[Configuration-selected coordinates]\label{cor:random-set}
Let $A=A(X)$ be an almost surely finite subset of $I$ selected measurably from
the sampled configuration.  Then
\begin{equation}\label{eq:random-set-cdf}
  \P(T_{A(X)}\le t\mid X)=F(t)^{|A(X)|}
  \quad\text{almost surely}.
\end{equation}
More generally, conditionally on $X$, the marks attached to the selected
coordinates remain independent $\GammaLaw(1/2,8)$ variables.
\end{corollary}

\begin{proof}
Condition on $X$ and apply the product law in Theorem~\ref{thm:main} to the
now deterministic finite set $A(X)$.
\end{proof}

\subsection{Exact confidence horizons}

\begin{corollary}[Exact and elementary horizons]\label{cor:horizons}
Fix $M\ge1$ and $0<\delta<1$.  The unique time $t_{M,\delta}$ satisfying
$\P(T_A\le t_{M,\delta})=1-\delta$ for $|A|=M$ is
\begin{equation}\label{eq:exact-horizon}
  t_{M,\delta}
  =4\left[
  \Phi^{-1}\!\left(
    \frac{1+(1-\delta)^{1/M}}{2}
  \right)\right]^2.
\end{equation}
The simpler sufficient condition
\begin{equation}\label{eq:simple-horizon}
  t\ge8\log(M/\delta)
\end{equation}
implies $\P(T_A\le t)\ge1-\delta$.
\end{corollary}

\begin{proof}
Equation \eqref{eq:exact-horizon} is obtained by solving $F(t)^M=1-\delta$.
The Gaussian Chernoff bound
$2\Phi(-x)\le e^{-x^2/2}$ for $x\ge0$ gives $q(t)\le e^{-t/8}$.
Therefore
\[
  \P(T_A>t)\le Mq(t)\le Me^{-t/8},
\]
which proves \eqref{eq:simple-horizon}.
\end{proof}
